\section{Principle and Action 15}

\paragraph{Chapter 15}

\begin{verse}
The Masters of ancient times\\
Were free, clairvoyant, mysterious, intuitive\\
In the vastness of the forces of their spirit they did not know of an I\\
This unconsciousness of the inner force gave greatness to their appearance.\\
To characterize them with images:\\
They were prudent, like those who traverse a stream of water in winter ice \\
Watchful, like those who know the enemy around themselves\\
Cold, like a stranger\\
Dissolving elusive , like ice that thaws\\
Rough, like rough-hewn wood\\
Vast, like the great valleys\\
Impenetrable, like turbid water.\\
Who today with the greatness of his light could lighten the inner darkness?\\
Who today with the greatness of his life could reanimate inner death inertia ?\\
In those men was Life\\
They were individuals and lords of the I\\
And their non-being their emptiness\footnote{non-grasping}, taking, or filled up in the figurative sense \\
In them was perfection.
\end{verse}

The characterization of the Masters of ancient times, models of Taoism, appear associated with the traditional idea of the regressive course of historical humanity, warned in a living way in the Far-Eastern world (a Chinese proverb says: ``Antiquity was like a full laughter, we are like empty shells'' --- see also the reference of Lao Tzu to the ``way of the Ancients'' in Chapter 14.) It is about a teaching of general importance, and whatever in the Tao Te Ching reflects it, must not be related to merely local Chinese political conditions, to the decadence of the Chou dynasty, under which Lao Tzu would have lived, like the banal interpretation of different commentators or translators.

\begin{quotex}
A living light endangered high antiquity, but some of its rays have reached us. To us it seems as if the Ancients lived in darkness because we see them through the thick fog from which we are emerging. ``Man is a child born at midnight: when he sees the sun rise, he believes that yesterday never existed'' (Taoist text cited by R. Remusat). See Chuang Tzu (XXXIII, 1) ``The ancients collaborated with the transcendent celestial and terrestrial influences, with the action of Heaven and Earth'' (it is the concept, previously noted, of the Great Triad)
\end{quotex}


Because of their structure in the original, the last lines of the chapter are controversial. It essentially followed Ular's reading as that which connects organically such lines with the rest and that maintains them at the level of the whole.

\paragraph{Chinese text and literal translation}

Chapter 15 (\begin{CJK*}{UTF8}{bsmi}第十五章\end{CJK*})

\columnratio{.4}
\begin{paracol}{2}
\begin{CJK*}{UTF8}{bsmi}
\begin{verse}
古之善為道者，微妙\\
玄達，深不可識。\\
夫唯不可識，故強為之容。\\
豫兮，若冬涉川；\\
猶兮，若畏四鄰；\\
儼兮，其若客；\\
渙兮，若冰之將釋；\\
敦兮，其若樸；\\
曠兮，其若谷；\\
渾兮，其若濁。\\
孰能濁以靜之徐清。\\
孰能安以動之徐生。\\
保此道不欲盈。\\
夫唯不盈，故能蔽而新成。
\end{verse}
\end{CJK*}
\switchcolumn
\begin{verse}
Once upon a time, those who knew the Way, were a mysterious and subtle people,\\
Transient yet profound, tranquil yet utterly unfathomable.\\
Since they are inexplicable, I can only tell what they seem like:\\
Cautious, as if wading through a winter river,\\
Wary, as if afraid of their own neighbors.\\
Grave, like the courteous house guests.\\
Elusive, as of melting ice.\\
Pure and natural, as of unchiseled gems.\\
Wide and open, as of a deep valley.\\
Yet mysterious, oh yes, they were like troubled water.\\
Who can remain tranquil amidst troubled airs, that calmness may flow from within?\\
Who may remain at peace eternal, that motion would yield birth to nature?\\
For those who follow the Way, fulfillment has never been their aim.\\
Only as they are forever unfulfilled, can such freshness be ever renewed.
\end{verse}
\end{paracol}

\flrightit{Posted on 2020-12-06 by Lao Tzu}

